\begin{filecontents*}{references.bib}
@book{lamport1994latex,
  author    = {Leslie Lamport},
  title     = {LaTeX: A Document Preparation System},
  year      = {1994},
  publisher = {Addison-Wesley}
}

@article{smith2024workflow,
  author  = {Smith, John},
  title   = {Editorial Workflows in Scientific Publishing},
  journal = {Journal of Documentation Models},
  year    = {2024},
  volume  = {18},
  number  = {2},
  pages   = {45-60}
}
\end{filecontents*}

\documentclass[conference]{IEEEtran}

\usepackage[brazil]{babel}
\usepackage[utf8]{inputenc}
\usepackage[T1]{fontenc}
\usepackage{cite}
\usepackage{graphicx}
\usepackage{amsmath,amssymb}
\usepackage{booktabs,multirow,array}
\usepackage[hidelinks]{hyperref}

\title{Modelo de Artigo Científico no Estilo IEEE}

\author{
\IEEEauthorblockN{Nome do Autor}
\IEEEauthorblockA{
Universidade / Instituição\\
email@exemplo.com
}
}

\begin{document}
\maketitle

\begin{abstract}
Este modelo apresenta a estrutura básica de um artigo científico no estilo IEEE, incluindo resumo, palavras-chave, seções com referências cruzadas, figura com legenda e bibliografia externa gerenciada via BibTeX.
\end{abstract}

\begin{IEEEkeywords}
LaTeX, IEEE, BibTeX, editoração científica, documentação técnica
\end{IEEEkeywords}

\section{Introdução}
A padronização de artigos científicos é essencial para garantir leitura fluida, consistência visual e facilidade na gestão editorial. O estilo IEEE é amplamente utilizado em áreas de engenharia e computação \cite{lamport1994latex}.

\section{Metodologia}
\label{sec:metodologia}
Como mostrado na Seção~\ref{sec:resultados}, a organização do texto e o uso adequado de referências cruzadas contribuem para a clareza do manuscrito. O fluxo editorial também pode ser tratado de forma estruturada \cite{smith2024workflow}.

\subsection{Referência cruzada}
O conteúdo desta seção será retomado na discussão apresentada na Seção~\ref{sec:resultados}.

\section{Resultados}
\label{sec:resultados}
Os resultados podem ser descritos em conjunto com gráficos, tabelas e análises comparativas. A Figura~\ref{fig:placeholder} ilustra um placeholder para inserção de imagem.

\begin{figure}[!t]
    \centering
    \fbox{\rule{0pt}{4cm}\rule{0.95\linewidth}{0pt}}
    \caption{Figura de exemplo com placeholder}
    \label{fig:placeholder}
\end{figure}

\begin{table}[!t]
\centering
\caption{Exemplo de tabela no estilo IEEE}
\label{tab:exemplo}
\begin{tabular}{lcc}
\toprule
Item & Valor A & Valor B \\
\midrule
Linha 1 & 10 & 12 \\
Linha 2 & 15 & 18 \\
Linha 3 & 20 & 25 \\
\bottomrule
\end{tabular}
\end{table}

\section{Discussão}
A leitura do material evidencia a importância de organização textual, consistência na formatação e uso correto de bibliografia externa. Isso facilita a revisão por pares e a publicação final.

\section{Conclusão}
Este template foi montado para demonstrar domínio da estrutura IEEE, incluindo abstração, seções numeradas, figura, tabela, citações e BibTeX.

\bibliographystyle{IEEEtran}
\bibliography{references}

\end{document}